\section{Execution Loop and Worked Example}\label{app:algorithm}

Algorithm~\ref{alg:cers} summarizes one execution of
\ours{} under the candidate budget $K$ and the review
budget $L$ for a single verification node.

\begin{algorithm}[t]
\caption{\ours{} execution loop}
\label{alg:cers}
\begin{algorithmic}[1]
\State inputs: task $q$; base agent $\mathcal{F}_\theta$
and the original action interface; chain budget $K$;
review budget $L$.
\State initialize rule state $M \leftarrow \emptyset$;
execution history $H \leftarrow \emptyset$.
\For{each step $t$}
  \State receive observation $o_t$; update chart
  references and interpretation context $c_t$.
  \State read applicable rules, revision evidence, and
  disputes from $M$.
  \State $(\widetilde{a}_t, \widetilde{\pi}_t)
  \leftarrow \mathcal{F}_\theta(q, o_t, H,
  \mathcal{R}_t)$.
  \If{a trigger holds: first chart-dependent judgment,
  interpretation conflict, or applicability change}
    \State expand the current explanation into
    $O \land B \Rightarrow C_q(\widetilde{a}_t)$.
    \State generate at most $K$ substantially different
    competing chains on the same complete chart.
    \State normalize candidates; merge repeated chains
    and shared rules.
    \State verify observations, rule applicability, and
    derivation sufficiency.
    \State apply adoption, revocation, replacement, or
    scope adjustment by the evidence.
    \State re-derive the current action from the updated
    state; within $L$, resolve remaining key disputes.
  \EndIf
  \State if the proposal reuses a revoked rule without
  new evidence, return the revision record and require
  re-derivation; with an effective new challenge,
  re-verify the rule.
  \State with \textsc{Keep} or \textsc{Revise}, execute
  the approved action through the original interface;
  with exhausted budget, record \textsc{Unresolved} and
  stop the decision.
  \State append the true observation, proposal, rule
  versions, action, and feedback to $H$.
\EndFor
\State save the task's rule revisions and execution
trajectory; close the task's rule scope.
\end{algorithmic}
\end{algorithm}

\subsection{Rule Merging by Object and Scope}\label{app:rule-merging}

``Darker means higher risk'' and ``a deeper color stands
for a more severe risk'' map to one rule when object and
scope coincide, whereas comparing Illinois with Kansas
and Illinois with Delaware under the same legend are two
comparisons that shared rules do not merge.


\subsection{Update Objects and Chain Exclusion}\label{app:update-objects}

When a chain is challenged, the method distinguishes
three objects of update: an \emph{observation error} (a
misread number, wrong entity binding, or wrong range)
revises the affected $O$ and re-verifies dependent
interpretations, never charging unreliable observations
to $B$; an \emph{inapplicable rule}---whose observation
stands while its observation-to-semantics mapping
disagrees with the current chart---is revoked or
scope-narrowed; and \emph{insufficient derivation}
leaves both supported while the task conclusion remains
unproven, so the module adds the missing comparison or
keeps an undetermined status, as when a higher candidate
refutes ``the current candidate is highest'' without
confirming the alternative is globally highest.

Conflicting conclusions between candidate chains
establish only that a conflict exists.
A chain is excluded only when chart-supported evidence
negates its observation, its rule applicability, or its
derivation conditions.

\subsection{A Worked Example of Revision and Cross-Step Reuse}\label{app:example}
\label{sec:method:example}

We illustrate the mechanism on the chart semantics of
task pub013 from our execution data, which asks the
agent to choose the highest-risk record among Illinois,
Kansas, and Delaware; the case material records that the
map legend assigns lighter yellow to higher risk and
deep red to lower risk.
The walkthrough below demonstrates the method and does
not add experimental results.
Suppose the current proposal selects Kansas under
\begin{equation}\label{eq:pub013}
\pi_0:
\underbrace{\text{Kansas is darker}}_{O_0}
\land
\underbrace{\text{darker means higher risk}}_{B_0}
\Rightarrow
C_q(\mathrm{KS}).
\end{equation}
The competing module asks why a darker color should mean
a higher risk; the model generates legend-dependent
interpretations on the same complete map and compares
the legend with the default color convention.
If verification confirms that the chart does not support
$B_0$, the rule loses applicability on this map and the
state records the original interpretation, the revision
reason that the convention contradicts the current
legend definition, the current interpretation that
candidate colors are read through the verified legend
mapping, the scope of the current map and risk metric,
and the condition for further updates, namely a
different encoding or an effective refutation of the
legend reading.
The agent then re-compares the candidates with $B_1$ and
the observed colors and completes the judgment the
public task requires; the concrete option is an output
of this reasoning, while the rule state keeps the
interpretation basis of the current map.
When the agent enters later pages or re-revisits the
chart-dependent choice, what it reads is that rule and
its revision reason: a new explanation that merely
repeats ``deep red usually means more danger'' must
answer the recorded legend evidence and cannot restore
$B_0$ directly, while a re-reading that finds the
earlier legend verification mistaken is allowed to
revise again.
The process records three distinct things: what changed
in the interpretation rules, whether the current action
must change, and which rule version later steps actually
used.
\section{Additional Results and Boundary Statements}\label{app:details}

\subsection{Matched-Configuration Seed Replication}
The revision experiments rerun direct selection and the
full configuration under three request seeds
(12345, 22345, 32345) on a single-GPU deployment whose
engine configuration matches the frozen service except
for batch admission (\texttt{max\_num\_seqs}=8 instead
of 1), so all comparisons reported for these runs are
internal to that configuration.
The three plain runs also provide the three samples for
the self-consistency baseline at no additional requests,
with majority vote among non-null selections and ties
recorded as abstentions.
The tabular-answering baseline re-requests each task in
two stages under the same engine: a transcription
request that renders, from the image, a compact
markdown table limited to the series and x positions
the public task's comparison involves (at most 24
rows), and a decision request answered from the table
alone, so the decision stage never sees the image.
The clean-chart control runs plain direct selection on
the paired clean charts and is scored against the
original labels, the protocol established by the
companion benchmark's online evaluation.

\subsection{Development Subset and Stratified Results}
The 24-task development subset is scored separately in
every round.
In the original round the full configuration V5 gains
$+17$ paired decisions on the remaining 116 tasks and
$+1$ on the development subset; the same decomposition
holds for the 8B round ($+10$ splits as $+9$ and
$+1$).
The environment domain is the only one in which V5
loses against direct selection in both model rounds
($-5$ in the 27B replication, $-9$ in the 8B round).
At the family level, subset-extrapolation tasks
contribute $+18$ (original round) and $+11$ (8B round)
for V5; dual-axis and non-linear-axis families are
mildly negative for V5 in the second model round
($-5$ and $-6$), and bar-height tasks are near-neutral
($-3$ and $-2$).

\subsection{Interface-Failure Accounting}
Interface failures are runs in which the agent never
produces a schema-valid selection.
They are zero for Direct, V4, V5, V6, and V7 on the
27B model (all three rounds), and 6, 8, and 51 for V3
in the original, replication, and 8B rounds
respectively.
V3's interface failures are not discarded: they remain
in the denominator, which is why its large
error-to-target correction counts on the 27B model
(30 forward against 6 reverse, exact McNemar
$p=0.0001$) reduce to a net $+2$.
The 8B actor's failure mode is dominated by responses
that violate the typed-record contract of the
verification pass rather than by timeouts or
truncations.

\subsection{Online Deployment Details}
Each online unit executes one task on one chart variant
under one system inside a sealed environment: public
inputs are mounted read-only, rule hooks intercept
chart evidence entering the context, and the unit ends
in exactly one of the recorded terminal statuses
(submitted, abstention with reason, re-verification
budget exhausted, or engineering failure).
The strict interface (\texttt{typed\_scalars\_v2})
permits at most eight actor calls and one
re-verification round per unit at temperature zero.
The three actor systems differ only below the
interface: a local 27B model under strict response
formats (Batch A), the same model behind a
schema-guided proxy mirroring the sealed structural
rules (Batch B), and the frontier model reached
directly through the gateway (Batch C and its
stability rerun).
Across the two frontier-model batches the defense arm
read its persistent rule state 48 and 46 times across
its 24 units, never cooccurred with a wrong
submission, and every per-unit persistence-causality
flag reads \texttt{not\_established}; the ordinary arm
produced exactly six wrong submissions, all
chart-induced, all under the 27B actor on misleading
charts, and none under the frontier actor.

\subsection{Boundary Statements}
Three boundary statements delimit the claims.
First, the per-unit \texttt{causal\_persistence}
effect is recorded as \emph{not established} in every
online unit: the harness shows that revised rules were
re-read and that outcomes changed between arms, but a
single-unit causal attribution of outcome change to
persistence is not identified by this design.
Second, the stability rerun's unit-level agreement
(37/48) bounds the reliability of any single-batch
online estimate; disaggregated four-way cells shift by
at most two units, and no flip ever crosses into a
wrong submission.
Third, the online deployment's paired effect is
one-directional: the defense arm converts ordinary
decisions into abstentions (10 and 13 per batch) and
never the reverse, so the deployment should be read as
a safe-fail characterization of the defense under the
strict protocol, not as an accuracy improvement
result.
